% ! TEX root = ../am2-18-19.tex

\chapter{Ecuații cu derivate parțiale și suprafețe de cîmp}

\section{Ecuații cu derivate parțiale de ordinul întîi}
\index{EDP}
\index{EDP!ordinul întîi}
\indent\indent Forma generală a unei ecuații cu derivate parțiale (EDP) de ordinul
întîi pentru funcția necunoscută $ u = u(x, y, z) $ este:
\[
  P(x, y, z) \dfrac{\partial u}{\partial x} + Q(x, y, z) \dfrac{\partial u}{\partial y} + %
  R(x, y, z) \dfrac{\partial u}{\partial z} = 0,
\]
unde $ P, Q, R : \RR^3 \to \RR $ sînt funcții de clasă (cel puțin) $\kal{C}^1$.

Pentru rezolvare, se scrie sistemul simetric asociat, care are forma generală
\[
  \dfrac{dx}{P} = \dfrac{dy}{Q} = \dfrac{dz}{R}
\]
și i se determină integralele prime. Dacă $ I_1 = C_1 $ și $ I_2 = C_2 $ sînt
integralele prime ale sistemului, atunci soluția finală a ecuației inițiale
este:
\[
  u(x, y, z) = \varphi(I_1, I_2),
\]
unde $ \varphi $ este o funcție oarecare de clasă (cel puțin) $\kal{C}^1$.

\begin{remark}\label{rk:not-dp}
  Pentru simplitate, vom mai folosi notațiile cunoscute pentru derivate parțiale,
  anume $ u_x = \dfrac{\partial u}{\partial x} $ etc.
\end{remark}

\begin{remark}\label{rk:int-indep}
  Metoda de rezolvare de mai sus cere implică și verificarea independenței celor
  două integrale prime, în sensul exemplificat mai jos.
\end{remark}

\textbf{Exemplu:} Rezolvăm ecuația cu derivate parțiale:
\[
  (4z - 5y)u_x + (5x - 3z)u_y + (3y - 4x)u_z = 0.
\]

\emph{Soluție:} Scriem sistemul simetric asociat, care este:
\[
  \dfrac{dx}{4z - 5y} = \dfrac{dy}{5x - 3z} = \dfrac{dz}{3y - 4x},
\]
sistem care este valabil pe domeniul:
\[
  D = \{(x, y, z \in \RR^3) \mid 4z \neq 5y, 5x \neq 3z, 3y \neq 4x \}.
\]

Pentru a obține integralele prime, amplificăm prima fracție cu 3, a doua cu 4 și
a treia cu 5 și obținem:
\[
  \dfrac{3dx}{12z - 15y} = \dfrac{4dy}{20x - 12z} = \dfrac{5dz}{15y - 20x} = %
  \dfrac{3dx + 4dy + 5dz}{0},
\]
de unde $ I_1(x, y, z) = 3x + 4y + 5z = c_1 $ este o integrală primă.

Mai putem amplifica și primul raport cu $ 2x $, pe al doilea cu $ 2y $ și
pe al treilea cu $ 2z $ și obținem:
\[
  \dfrac{2xdx}{8xz - 10xy} = \dfrac{2ydy}{10xy - 6yz} = \dfrac{2zdz}{6yz - 8xz} = %
  \dfrac{2xdx + 2ydy + 2zdz}{0},
\]
deci o a doua integrală primă este $ I_2(x, y, z) = x^2 + y^2 + z^2 = c_2 $.

Evident, cele două integrale prime sînt definite pe același domeniu $ D $.

Independența celor două înseamnă verificarea că matricea dată de derivatele
lor par\-ți\-a\-le are rangul maxim pe domeniul de definiție. Avem, așadar:
\[
  A = \begin{pmatrix}
    I_{1x} & I_{2x} \\
    I_{1y} & I_{2y} \\
    I_{1z} & I_{2z}
  \end{pmatrix} = \begin{pmatrix}
    3 & 2x \\
    4 & 2y \\
    5 & 2z
  \end{pmatrix},
\]
matrice care se poate verifica imediat că are rangul 2, adică maxim, pentru orice
$(x,y,z) \in D $. Deci integralele prime sînt independente și soluția finală a
ecuației este:
\[
  u(x, y, z) = \varphi(3x + 4y + 5z, x^2 + y^2 + z^2),
\]
unde $ \varphi $ este o funcție arbitrară de clasă $ \kal{C}^1 $.

%%%%%%%%%%%%%%%%%%%%%%%%%%%%%%%%%%%%%%%%%%%%%%%%%%%%%%%%%%%%%%%%%%%%%% 

\section{Suprafețe de cîmp}
\index{suprafețe!de cîmp}

\indent\indent Suprafețele de cîmp pentru un cîmp vectorial tridimensional
se obțin scriind o ecuație cu derivate parțiale asociată, căreia îi determinăm
soluția generală. Suprafața de cîmp este, atunci, dată de anularea soluției
generale, scrisă în forma implicită.

\textbf{Exemplu:} Fie cîmpul vectorial:
\[
  \vec{V} = xy^2 \vec{i} + x^2y \vec{j} + z(x^2 + y^2) \vec{k}.
\]
Determinăm suprafețele de cîmp.

\emph{Soluție:} Scriem ecuația cu derivate parțiale asociată cîmpului,
pentru o funcție necunoscută $ u = u(x, y, z) $. Avem:
\[
  xy^2 u_x + x^2y u_y + z(x^2 + y^2)u_z = 0. 
\]
Rezolvarea ecuației se face ca mai sus, scriind sistemul asociat:
\[
  \dfrac{dx}{xy^2} = \dfrac{dy}{x^2y} = \dfrac{dz}{z(x^2 + y^2)},
\]
definit pe domeniul potrivit, adică $ D = \RR^3 - \{(0, 0, 0)\}$.

Din prima egalitate, putem simplifica cu $ xy $ și avem $ xdx = ydy $, deci
o integrală primă este $ x^2 - y^2 = c_1 $.

Putem amplifica rapoartele cu $ y, x $ și respectiv $ 1 $ și avem:
\[
  \dfrac{ydx}{xy^3} = \dfrac{xdy}{x^3y} = \dfrac{dz}{z(x^2 + y^2)} = %
  \dfrac{ydx + xdy}{xy(x^2 + y^2)}.
\]
Din ultimele două rapoarte, obținem, după simplificare cu $ x^2 + y^2 $:
\[
  \dfrac{dz}{z} = \dfrac{d(xy)}{xy},
\]
deci $ \dfrac{xy}{z} = c_2 $.

Acum soluția generală a ecuației cu derivate parțiale este:
\[
  u(x, y, z) = \varphi\Big(x^2 - y^2, \dfrac{xy}{z} \Big),
\]
cu $ \varphi $ o funcție arbitrară de clasă $ \kal{C}^1 $, deci suprafețele de
cîmp ale cîmpului $ \vec{V} $ au forma:
\[
  \varphi\Big( x^2 - y^2, \dfrac{xy}{z} \Big) = 0.
\]

\vspace{1cm}

În unele cazuri, se pot cere anume suprafețe de cîmp, de exemplu:

\textbf{Exemplu:} Să se determine suprafața de cîmp a cîmpului vectorial:
\[
  \vec{V} = yz\vec{i} + xz \vec{j} + xy\vec{k},
\]
care trece prin curba dată de intersecția cilindrului $ x^2 + z^2 = 4 $ cu
planul $ y = 0 $.

\emph{Soluție:} Mai întîi, determinăm toate suprafețele de cîmp, ca mai sus.
Scriem direct sistemul asociat:
\[
  \dfrac{dx}{yz} = \dfrac{dy}{xz} = \dfrac{dz}{xy}.
\]
Amplificăm cu $ x, y $, respectiv $ z $ și obținem $ xdx = ydy = zdz $, adică
$ x^2 - y^2 = c_1, x^2 - z^2 = c_2 $ sînt două integrale prime.

Aceste integrale prime dau o familie infinită de suprafețe, dintre care
vrem să aflăm pe cea care trece prin curba dată. Așadar, avem de rezolvat
sistemul de ecuații:
\[
  \begin{cases}
    x^2 - y^2 &= c_1  \\
    x^2 - z^2 &= c_2 \\
    x^2 + z^2 &= 4 \\
    y &= 0
  \end{cases}
\]
pentru constantele $ c_1 $ și $ c_2 $, care ne vor identifica exact suprafața.

Este suficient să găsim condiția de compatibilitate a sistemului, care se poate
obține din primele două ecuații. Înmulțim prima ecuație cu 2 și o scădem din ea
pe a doua și comparăm cu a treia ecuație. Se obține $ 2c_1 - c_2 = 4 $.
Apoi, înlocuim în integralele prime și avem $ 2(x^2 - y^2) - (x^2 - z^2) = 4 $, care
se prelucrează și se aduce la forma canonică:
\[
  \dfrac{x^2}{4} - \dfrac{y^2}{2} + \dfrac{z^2}{4} = 1,
\]
care este un \emph{hiperboloid cu o pînză}.

%%%%%%%%%%%%%%%%%%%%%%%%%%%%%%%%%%%%%%%%%%%%%%%%%%%%%%%%%%%%%%%%%%%%%% 

\section{Ecuații cu derivate parțiale cvasiliniare}
\index{EDP!cvasiliniare}

\indent\indent Forma generală a ecuațiilor cu derivate parțiale cvasiliniare
pentru o funcție $ u = u(x_1, x_2, \dots, x_n) $ este:
\[
  P_1(x_1, \dots, x_n) u_{x_1} + P_2(x_1, \dots, x_n)u_{x_2} + \dots + %
  P_{n-1}(x_1, \dots, x_n) u_{x_{n-1}} + Q(x_1, \dots, x_n) = 0,
\]
adică apar toate derivatele parțiale ale lui $ u $, mai puțin ultima.

În exercițiile pe care o să le întîlnim cel mai des, funcția
$ u $ este înlocuită cu $ z = z(x, y) $, iar atunci ultima derivată
parțială o putem gîndi, de exemplu, ca $ z_z = 1 $.

Modul de rezolvare este explicat pe exemplul de mai jos.

\textbf{Exemplu:} Fie ecuația cvasiliniară:
\[
  (z - y)^2 \dfrac{\partial z}{\partial x} + xz \dfrac{\partial z}{\partial y} = xy.
\]
\emph{Soluție:} Se poate vedea că ecuația este cvasiliniară, funcția
necunoscută fiind $ z = z(x, y) $.

Din teorie, ecuația trebuie să aibă o soluție scrisă în formă implicită
$ u(x, y, z) = 0 $. De fapt, această funcție trebuie înțeleasă ca
$ u = u(x, y, z(x, y)) $. Folosind formula de derivare a funcțiilor compuse,
dacă vrem să scriem derivata parțială în raport cu $ x $ a lui $ u $, trebuie
să ținem cont că $ x $ apare și în $ z(x, y) $, deci avem:
\[
  0 = \dfrac{\partial u}{\partial x} + \dfrac{\partial u}{\partial z} \cdot %
  \dfrac{\partial z}{\partial x},
\]
și similar pentru derivata în raport cu $ y $. Rezultă:
\[
  z_x = -\dfrac{u_x}{u_z}, \quad z_y = -\dfrac{u_y}{u_z}.
\]
Înlocuim în ecuația dată și înmulțim relația cu $ u_z $, obținînd:
\[
  (z - y)^2 u_x + xz u_y + xy u_z = 0,
\]
care este o ecuație cu derivate parțiale de ordinul întîi și o putem rezolva
ca în prima secțiune.

Scriem sistemul asociat:
\[
  \dfrac{dx}{(z - y)^2} = \dfrac{dy}{xz} = \dfrac{dz}{xy}.
\]
Din a doua egalitate obținem direct $ y^2 - z^2 = c_1 $, care este o integrală primă.

Amplificăm primul raport cu $ x $, pe al doilea cu $ (y - z) $ și pe al treilea cu $ (z - y)$ 
și obținem prin adunare:
\[
  xdx + (y - z)dy + (z - y) dz = 0 \Rightarrow xdx + ydy + zdz - d(yz) = 0.
\]
Rezultă a doua integrală primă $ x^2 + y^2 + z^2 - 2yz = x^2 + (y - z)^2 = c_2 $.

Soluția pentru $ u $ va fi:
\[
  u(x, y, z) = \varphi(y^2 - z^2, x^2 + (y - z)^2),
\]
cu $ \varphi $ o funcție de clasă $ \kal{C}^1 $ pe domeniul de definiție.

Atunci soluția pentru $ z $ se obține din forma implicită $ u(x, y, z) = 0 $.


%%%%%%%%%%%%%%%%%%%%%%%%%%%%%%%%%%%%%%%%%%%%%%%%%%%%%%%%%%%%%%%%%%%%%% 

\section{Exerciții}

\indent\indent 1. Rezolvați ecuațiile cvasiliniare:
\begin{enumerate}[(a)]
\item $x(y^3 - 2x^3)z_x + y(2y^3 - x^3)z_y = 9z(x^3 - y^3) $;
\item $ 2xzz_x + 2yzz_y = z^2 - x^2 - y^2 $;
\item $ 2y z_x + 3x^2 z_y + 6x^2 y = 0 $;
\item $ x \Big( y^2 - z^2 \Big) z_x - y \Big( x^2 + y^2 \Big) z_y = z \Big( x^2 + y^2 \Big) $;
\item $ xz z_x + yz z_y = x + y $.
\end{enumerate}

\emph{Indicații:} (a) Sistemul diferențial autonom la care se ajunge este:
\[
  \dfrac{dx}{x(y^3 - 2x^3)} = \dfrac{dy}{y(2y^3 - x^3)} = \dfrac{dz}{9z(x^3 - y^3)}.
\]

Rezultă:
\[
  \dfrac{ydx + xdy}{3xy(y^3 - x^3)} = \dfrac{dz}{9z(x^3 - y^3)} \Rightarrow %
  \dfrac{d(xy)}{xy} = -\dfrac{dz}{3z} \Rightarrow x^3 y^3 z = c_1.
\]

Din primele două rapoarte obținem:
\[
  \dfrac{y(2y^3 - x^3)}{x(y^3 - 2x^3)} = \dfrac{dy}{dx}.
\]

Simplificăm forțat cu $ \dfrac{y}{x} = u $ și ajungem la ecuația:
\[
  \dfrac{u^3 - 2}{u(u+1)(u^2 - u + 1)} du = \dfrac{dx}{x}.
\]

Rezultă $ \ln(c_2 x) = -2\ln u + \ln(u + 1) + \ln(u^2 - u + 1) $, adică, în final,
$ c_2 = \dfrac{x^3 + y^3}{x^2 y^2} $.

(b) Se ajunge la sistemul autonom:
\[
  \dfrac{dx}{2xz} = \dfrac{dy}{2yz} = \dfrac{dz}{z^2 - x^2 - y^2}.
\]

Din primele două rapoarte avem $ \dfrac{y}{x} = c_1 $ și apoi:
\[
  \dfrac{2xdx}{2x^2} = \dfrac{2ydy}{2y^2} = \dfrac{2zdz}{z^2 - x^2 - y^2},
\]
adică $ c_2x = x^2 + y^2 + z^2 $.

\vspace{1cm}

2. Determinați suprafețele de cîmp pentru cîmpurile vectoriale:
\begin{enumerate}[(a)]
\item $ \vec{V} = (x + y + z)\vec{i} + (x - y)\vec{j} + (y - x)\vec{k} $;
\item $ \vec{V} = (x - y)\vec{i} + (x + y) \vec{j} + \vec{k} $.
\end{enumerate}

\emph{Indicații:} (a) Ajungem la sistemul:
\[
  \dfrac{dx}{x + y + z} = \dfrac{dy}{x - y} = \dfrac{dz}{y - x}.
\]

Din prima egalitate, rezultă $ dy + dz = 0 \Rightarrow y + z = c_1 $. Înlocuind
în cea de-a doua egalitate $ z = c_1 - y $, rezultă $ (x - y)dx = (x + c_1)dy $,
deci $ xdx - d(xy) - c_1 dy = 0 $, adică $ \dfrac{x^2}{2} - xy - c_1y = c_2 $. Înlocuim
$ c_1 = y + z $ și obtinem $ \dfrac{x^2}{2} - xy - y^2 - yz = c_2 $.

Domeniul de definiție va fi $ D = \{ (x, y, z) \in \RR^3 \mid x + y + z \neq 0 \} $.

(b) Se ajunge la sistemul:
\[
  \dfrac{dx}{x - y} = \dfrac{dy}{x + y} = \dfrac{dz}{1}.
\]

Rezultă:
\[
  \dfrac{xdx}{x^2 - xy} = \dfrac{ydy}{xy + y^2} = \dfrac{xdx + ydy}{x^2 + y^2} = %
  \dfrac{\frac{1}{2} d(x^2 + y^2)}{x^2 + y^2} = \dfrac{dz}{1},
\]
deci $ x^2 + y^2 = c_1 e^{2z} $, iar apoi, din $ \dfrac{dy}{dx} = \dfrac{x + y}{x - y} $,
putem nota $ \dfrac{y}{x} = u $ și rezultă $ x^2 + y^2 = c_2 e^{2 \arctan \frac{y}{x}} $.

\vspace{1cm}

3. Determinați suprafața de cîmp a cîmpului vectorial:
\[
  \vec{V} = 2xz\vec{i} + 2yz\vec{j} + (z^2 - x^2 - y^2)\vec{k},
\]
care conține cercul dat de $ z = 0 $ și $ x^2 + y^2 - 2x = 0 $.

\emph{Indicație:} Integralele prime independente ale sistemului se
determină din:
\[
  \dfrac{dx}{2xz} = \dfrac{dy}{2yz} = \dfrac{dz}{z^2 - x^2 - y^2}.
\]
Din primele două rapoarte, obținem $ y = c_1 x $, iar apoi, putem înmulți
toate egalitățile cu $ 2z $ și prelucrăm mai departe:
\[
  \dfrac{dx}{x} = \dfrac{dy}{y} = \dfrac{2z}{z^2 - x^2 - y^2} \Rightarrow %
  \dfrac{2xdx}{2x^2} = \dfrac{2ydy}{2y^2} = \dfrac{2zdz}{z^2 - x^2 - y^2} = %
  \dfrac{2(xdx + ydy + zdz)}{x^2 + y^2 + z^2},
\]
deci $ x^2 + y^2 + z^2 = c_2 x $.

Pentru a obține intersecția cu cercul dat, avem condiția ca sistemul de mai jos să fie
compatibil:
\[
  \begin{cases}
    y &= c_1 x \\
    x^2 + y^2 + z^2 &= c_2 x \\
    z &= 0 \\
    x^2 + y^2 - 2x &= 0
  \end{cases}.
\]

Din ultimele trei ecuații se obține că sistemul este compatibil dacă și numai dacă $ c_2 = 2 $.

Rezultă $ x^2 + y^2 + z^2 - 2x = 0 $, care este o sferă.

\vspace{1cm}

4. Fie cîmpul vectorial:
\[
  \vec{V} = (x + y)\vec{i} + (y - x)\vec{j} - 2z\vec{k}.
\]
Să se determine:
\begin{enumerate}[(a)]
\item liniile de cîmp;
\item linia de cîmp ce conține punctul $ M(1, 0, 1) $;
\item suprafețele de cîmp;
\item suprafața de cîmp care conține dreapta $ z = 1, y - x\sqrt{3} = 0 $.
\end{enumerate}

\emph{Indicație:} Pentru liniile de cîmp, ecuația dată de primele două rapoarte:
\[
  \dfrac{dx}{dy} = \dfrac{x + y}{y - x}
\]
este o ecuație diferențială de ordinul întîi, omogenă, care se poate rezolva cu
substituția $ y = tx $.

\vspace{1cm}

5. Să se determine soluția ecuațiilor cu derivate parțiale de ordinul întîi:
\begin{enumerate}[(a)]
\item $ y u_x - x u_y = 0 $;
\item $ xz u_x - yz u_y + (x^2 + y^2) u_z = 0 $;
\item $ x u_x - y u_y = 0 $.
\end{enumerate}

%%% Local Variables:
%%% mode: latex
%%% TeX-master: "../am2-18-19"
%%% End:
